% Registro de lo recortado de cap3/3_4_fallas.tex (no se carga con \input).

% ===================================================================
% RECORTADO de cap3/3_4_fallas.tex: §4.4.9 El punto perturbado
% ===================================================================
\subsection{Una comprobación del registro: el punto perturbado}
\label{sec:registro_corrige}

Durante la sesión del 8 de septiembre de 2026 hubo una interrupción
física del arreglo mientras la secuencia estaba en curso: se colocó una
protección sobre el láser con la adquisición corriendo. El hecho quedó
anotado en la bitácora de la sesión junto con la atribución que el
operador hizo en ese momento, señalando dos capturas concretas como
afectadas.

El análisis posterior de los archivos, realizado días después y
sin más elementos que los datos guardados, contradice esa atribución.
Las dos capturas señaladas no se distinguen de sus vecinas por ningún
criterio: ni su amplitud pico a pico, que sigue la tendencia
descendente de la serie, ni el instante de su máximo, que coincide con
el del resto dentro de medio microsegundo.

La captura que sí se aparta es \texttt{m0004}. Dos indicadores
independientes la señalan. Su amplitud pico a pico es
\SI{1.732}{\milli\volt}, el valor más alto de toda la sesión, en medio
de una serie que desciende de forma sostenida. Y el instante de su
máximo absoluto cae en \SI{78.08}{\micro\second}, cuando en las ocho
capturas restantes se ubica entre \SI{94.37}{\micro\second} y
\SI{94.89}{\micro\second}. No se trata de la misma señal medida con
ruido: es una señal distinta.

El episodio ilustra el argumento de este capítulo mejor que ninguna de
las anomalías anteriores, y conviene enunciarlo con precisión para no
reclamar de más. El sistema no detectó la perturbación: ningún campo
del archivo levantó bandera, y no debía hacerlo, porque nada en la
adquisición había fallado. Lo que el sistema hizo fue registrar, junto
a cada forma de onda, las condiciones bajo las que se obtuvo y la
trazabilidad suficiente para que la pregunta pudiera formularse
después. La perturbación se identificó comparando nueve capturas entre
sí, algo que solo es posible si las nueve están guardadas con sus
escalas, su temperatura y su orden.

Ese es el criterio con el que debe juzgarse un sistema de
automatización de medición: no que anticipe todo lo que puede salir
mal, sino que deje registro suficiente para reconstruirlo. En este
caso el registro corrigió al operador, que estuvo presente, sobre lo
que el operador creyó haber visto.

